
\examquestion{2023}{3}{Squeezed-coherent input in a Mach--Zehnder interferometer}{4}
We have one squeezed-coherent input mode \(a\) and
one vacuum input mode \(v\).
For clarity, distinguish \emph{field displacement}
\(\alpha\) from squeeze parameters
\(\mu=\cosh r,\ \nu=e^{2i\varphi_S}\sinh r\).

\subsection*{(a) Energy conservation at the first beam splitter}
The stated beam splitter matrix is
\[
M=\frac1{\sqrt2}\begin{pmatrix}1&i\\i&1\end{pmatrix}.
\]
The input column of field amplitudes is
\(\mathbf E_{\rm in}=(E_{\rm in},E_v)^T\).
Since
\[
M^\dagger M=\frac12
\begin{pmatrix}1&-i\\-i&1\end{pmatrix}
\begin{pmatrix}1&i\\i&1\end{pmatrix}=I,
\]
we obtain
\[
\boxed{|E_1|^2+|E_2|^2=
\mathbf E_{\rm in}^\dagger M^\dagger M
\mathbf E_{\rm in}
=|E_{\rm in}|^2+|E_v|^2}.
\]
For quantum modes the analogous statement is
conservation of total photon number,
because \(M\) is unitary.

\subsection*{(b) Will squeezing enhance the sensitivity?}
Not \emph{automatically}. Squeezing can lower
noise in a selected quadrature, but the effect
depends on which input port is squeezed, its
squeezing angle, the operating phase and the
measurement. In this question the \emph{other input
is vacuum} and the measured observable is
the \emph{output photon-number difference}.
As the calculations below show, at \(\phi=\pi/2\)
the variance of the difference is the
\emph{mean photon number} of the occupied input.
Thus this exact configuration does \emph{not}
give sub-shot-noise phase resolution at fixed
total mean photon number through simple
difference detection. Injecting appropriately
squeezed vacuum in the unused port is a different
well-known enhancement scheme.

\subsection*{(c) Derive the number-difference operator}
The problem gives
\[
a_A=c\,a+d\,v,\qquad a_B=d\,a+c\,v,\quad
c=\tfrac12(e^{i\phi}+1),\quad
d=\tfrac12(e^{i\phi}-1).
\]
Compute the coefficients:
\[
|c|^2-|d|^2=\cos\phi,\qquad
c^*d-d^*c=i\sin\phi.
\]
Expand \(a_A^\dagger a_A-a_B^\dagger a_B\):
\[
\boxed{D\equiv n_A-n_B
=\cos\phi(n_a-n_v)
+i\sin\phi(a^\dagger v-v^\dagger a)}.
\]

\subsection*{(d) Signal for the squeezed-coherent input}
Since the second mode is vacuum,
\(\langle n_v\rangle=0\),
\(\langle v\rangle=0\); all cross-mode
expectation values factor. Hence
\[
\boxed{\langle D\rangle=\cos\phi\,N},
\qquad N=\langle a^\dagger a\rangle.
\]
For the conventional displaced-squeezed
state \(\ket{\alpha,\xi}=D(\alpha)S(\xi)\ket0\),
we have
\[
S^\dagger a S=\mu a-\nu a^\dagger,\qquad
\boxed{N=|\alpha|^2+|\nu|^2
=|\alpha|^2+\sinh^2 r}.
\]
Therefore
\[
\boxed{\langle n_A-n_B\rangle
=\cos\phi(|\alpha|^2+|\nu|^2)}.
\]
If the problem instead defines
\(\ket{\alpha,\xi}=S(\xi)D(\alpha)\ket0\),
operator ordering changes the number expectation:
\(N=|\mu\alpha-\nu\alpha^*|^2+|\nu|^2\).
The subsequent noise formula in terms of \(N\)
remains valid. \emph{State the convention.}

\subsection*{(e) Number-difference variance at \(\phi=\pi/2\)}
At the operating phase,
\(D=i(a^\dagger v-v^\dagger a)\) and
\(\langle D\rangle=0\). Square it:
\[
D^2=-a^\dagger v a^\dagger v
+a^\dagger v v^\dagger a
+v^\dagger a a^\dagger v
-v^\dagger a v^\dagger a.
\]
In the vacuum \(v\),
\(\langle v^\dagger v\rangle=0\),
\(\langle vv^\dagger\rangle=1\),
\(\langle v^2\rangle=\langle v^{\dagger2}\rangle=0\).
Only the second term remains, so
\[
\boxed{\operatorname{Var}(D)|_{\pi/2}
=\langle D^2\rangle
=\langle a^\dagger a\rangle=N
=|\alpha|^2+|\nu|^2}
\]
for \(D(\alpha)S(\xi)\ket0\).

\subsection*{(f) Minimum resolvable phase near \(\pi/2\)}
For a small deviation \(\epsilon\),
\(\cos(\pi/2+\epsilon)\simeq-\epsilon\).
Thus the signal magnitude is
\(|\langle D\rangle|\simeq N|\epsilon|\),
and the noise standard deviation is \(\sqrt N\).
Demanding signal-to-noise ratio 1 gives
\[
1\simeq\frac{N|\epsilon|}{\sqrt N},
\qquad
\boxed{\delta\phi_{\rm min}\simeq\frac1{\sqrt N}
=\frac1{\sqrt{|\alpha|^2+\sinh^2 r}}}.
\]
For \(M\) independent repetitions,
the standard error scales as \(1/\sqrt{MN}\).
Squeezing can raise \(N\) if \(\alpha\) is
held fixed, but \emph{at equal total photon
number} this is still shot-noise scaling.
\solcheck{The sensitivity must be computed using
the \emph{slope} at the dark operating point:
the mean signal itself is zero exactly at
\(\phi=\pi/2\). The derivation does not assume
a Gaussian input; only that the other port is vacuum.}
